\chapter[La Red Vacía]{La Red Vacía\\ \large Quien saquea sin medida las aguas del río siembra hambre y engaño en su propia boca.}

\elegantimage{22_pesca_excesiva/web/assets/hero.jpg}

\lettrine[lines=3, nindent=0.5em, findent=0.2em]{E}{}ntre las transgresiones que el libro celeste registra con especial gravedad se encuentra la pesca excesiva y la explotación desmedida de los ríos y mares. Quien arranca de las aguas más de lo que necesita para vivir, quien vacía los estanques con redes insaciables sin dejar sustento para las criaturas que vendrán después, está rompiendo el pacto sagrado que une al hombre con la naturaleza que lo alimenta...

Cada pez arrancado del agua más allá de la necesidad legítima, cada red tendida con la avaricia como único motor, cada ecosistema acuático devastado por la codicia sin freno, deja una marca indeleble en el registro cósmico del alma. La deuda se acumula como agua turbia que un día desbordará.

La ley universal dicta con severidad incontestable que quien saquea los ríos y mares sin compasión por las generaciones futuras renacerá con la boca lacerada por llagas, sufriendo hambre crónica y siendo víctima constante del engaño, experimentando en su propio cuerpo la escasez que impuso a las aguas que saqueó.

\section{El Pescador Insaciable}
\begin{elegantquote}
\textit{En un delta fértil donde tres ríos confluían, vivía un pescador llamado Khanh cuya habilidad con las redes no tenía rival. Pero Khanh no pescaba para alimentar a su familia: pescaba para dominar. Tendía trampas desde el amanecer hasta la medianoche, arrasando bancos enteros de peces que luego dejaba pudrir al sol porque no podía vender ni consumir semejante abundancia...}

\textit{Los ancianos del delta le suplicaban moderación: "Deja algo para que los peces se reproduzcan, deja algo para quienes vendrán después de ti." Pero Khanh se reía de sus advertencias, convencido de que el río era inagotable y que su dominio sobre las aguas era un derecho ganado por su destreza.}

\textit{En su siguiente existencia, Khanh nació con una boca frágil plagada de llagas que le impedían comer sin dolor. Rodeado de abundante comida que no podía tragar, era además engañado constantemente por quienes prometían curarle. Miraba el río —ahora rebosante de peces que no podía saborear— y comprendía, entre lágrimas silenciosas, que cada bocado que le fue negado era un eco exacto de cada pez que él había arrebatado al agua sin necesidad.}

\end{elegantquote}

\elegantimage{22_pesca_excesiva/web/assets/art.jpg}

\vspace{1em}
\begin{center}
    \textbf{\Large Quien vacía las aguas sin compasión llena su propia boca de llagas y su vida de engaños.}
\end{center}
\vspace{1em}

\section{Análisis Kármico y Traducción}
\begin{wrapfigure}{r}{0.5\textwidth}
  \vspace{-1.5em}
  \begin{center}
  \inlineelegantimage[0.45\textwidth]{22_pesca_excesiva/web/assets/pasaje_original.jpg}
  \end{center}
  \vspace{-1.5em}
\end{wrapfigure}
La exhibición impúdica del cuerpo no es un acto inocente ni meramente estético: constituye, según la ley kármica, una forma de agresión sutil contra la paz interior de quienes contemplan. Al provocar deliberadamente el deseo ajeno, el exhibicionista siembra en el campo cósmico semillas de fuego que inevitablemente germinarán en su propia alma. La retribución es de una simetría perfecta: quien encendió el deseo de otros sin medida, renace consumido por un fuego interno que ninguna relación, ningún placer y ninguna satisfacción pueden apagar. Es la sed del alma que crece cuanto más se bebe.

